\section{Data and preprocessing}
\label{app:data}

\paragraph{CIFAR-Aug.} From the CIFAR-10 training set we draw $200$ images per class with a fixed seed. Each source receives $16$ variants per strength; variant $0$ is the untouched image. \emph{Weak}: brightness, contrast and saturation each scaled by a factor in $[0.85,1.15]$, then a horizontal flip with probability $0.5$. \emph{Medium}: flip ($p=0.5$), rotation uniform in $[-20^\circ,20^\circ]$, a square crop covering $75$--$100\%$ of the side at a random position resized back to $32\times32$, then colour factors in $[0.7,1.3]$. \emph{Strong}: as medium with rotation up to $30^\circ$, crops down to $60\%$ and colour factors in $[0.6,1.4]$, followed by one or two filters drawn from grayscale, Gaussian blur, sharpen, equalise, solarise, posterise, edge enhancement, quarter-turn, vertical flip and a per-channel intensity shift. The \emph{fresh} pool holds $300$ images per class from the CIFAR-10 test set. Generation is deterministic (seeded per source and strength).

\paragraph{Leaf-A and Leaf-B.} Two public releases of field photographs of crop leaves. Leaf-A has $9{,}000$ images: $600$ source photographs in three classes, each released with $15$ variants named \texttt{v1}--\texttt{v15}; \texttt{v1} appears to be the unedited photograph and the others include colour and contrast changes, grayscale, thresholding, edge maps, rotation and flips. Leaf-B has $13{,}600$ images: $850$ source photographs in three classes, each with $16$ named variants (Original, Flip, Rotate, Brightness, Contrast, Gamma, HSV, CLAHE, Grayscale, Gaussian and median blur, bilateral filter, sharpen, unsharp mask, top-hat, black-hat). In both releases the source identifier is part of the file name, which is what gives us ground-truth groups. To preserve double-blind review we do not name the releases here; they will be cited in the final version.

\paragraph{PlantDoc.} The official train ($2{,}342$) and test ($236$) folders of the classification release~\citep{singh2020plantdoc}: $28$ classes, $27$ of which occur in the test folder.

\paragraph{COVID-CXR.} From the metadata of the public collection~\citep{cohen2020covid} we keep X-ray images in the \texttt{images} folder with a frontal view (PA, AP, AP supine, AP erect) and a non-empty finding: $699$ images of $424$ patients. The label is COVID-19 versus any other finding; the group is the patient identifier.

\paragraph{CIFAR-10.} The official $50{,}000$/$10{,}000$ split.

\paragraph{Brain MRI.} The official \texttt{Training} ($2{,}870$) and \texttt{Testing} ($394$) folders of the four-class collection of \citet{bhuvaji2020brain} (glioma, meningioma, pituitary tumour, no tumour). Probes use frozen MobileNetV3 features at $160$ pixels.

\paragraph{Fruits-360.} The official \texttt{Training} and \texttt{Test} folders~\citep{muresan2018fruits} for $15$ classes (apple Braeburn, avocado, banana, cherry, kiwi, lemon, mango, orange, peach, pear, pineapple, pomegranate, strawberry, tomato, walnut): $7{,}759$ training and $2{,}600$ test images. Probes use frozen MobileNetV3 features at $112$ pixels.

\section{Hyperparameters}
\label{app:hyper}

\paragraph{Features.} timm~\citep{wightman2019timm} MobileNetV3-Large (\texttt{mobilenetv3\_large\_100}) and EfficientNet-B0 with public ImageNet weights, classifier removed (global-pooled features). Inputs are resized (no crop) to $112$ (CIFAR), $128$ (leaf releases) or $160$ pixels (PlantDoc, COVID-CXR) and normalised with the model's statistics. The dihedral orbit average is the mean embedding over the eight dihedral transforms.

\paragraph{Probes.} Linear probe: multinomial logistic regression on standardised features, $\ell_2$ penalty $10^{-3}$ on the weights, full-batch L-BFGS (50 iterations, strong Wolfe line search). 1-NN and 5-NN use cosine similarity. Splits are 80/20; group splits assign whole groups at random to reach a 20\% share of groups.

\paragraph{AugLeak.} Grayscale thumbnails $64\times64$, Gaussian blur $\sigma=1$, resampled to $32\times32$, disc of radius $0.48$ of the side ($732$ pixels). Query transforms: $8$ dihedral elements; rotations by $\pm15^\circ,\pm30^\circ,\pm45^\circ$ of the image and its mirror; centre crops of $80\%$ and $65\%$ of the side of the image and its mirror; corner crops of $75\%$ ($28$ transforms). Hash view: the $64$-bit DCT perceptual hash of the $32\times32$ thumbnail, searched over the $8$ dihedral elements. Hubness statistics use all other-class images. The Gumbel null is fitted by maximum likelihood to the best other-class score of every image. Views are combined by the Bonferroni-adjusted minimum $p$-value. Byte-identical same-class files are linked regardless of score (MD5 of the file). The audits of Section~\ref{sec:audit} use the structure view alone with this rule; the released command-line tool also adds the hash view and so flags somewhat more (for the brain-MRI split, $250$ rather than $200$ test images at $\alpha=10^{-2}$).

\paragraph{Calibration and the number of classes.} An earlier version of our CIFAR-Aug benchmark drew its $500$ sources from three classes. There, the pixel views linked same-class non-siblings at $1.5$--$3.1\times\alpha$ for $\alpha=10^{-2}$ and $10$--$60\times\alpha$ for $\alpha=10^{-4}$, against less than $\alpha$ on the ten-class subset. With few classes, the other-class best matches that define the null come from too narrow a slice of the image distribution. On releases with few classes, such as the two three-class leaf releases, the nominal $\alpha$ should therefore be read as a lower bound on the false-link rate.

\section{Invariant embedding}
\label{app:invariant}

We draw $M=1{,}500$ images from the release, make two copies of each with a generic edit family (horizontal and vertical flips, quarter-turns, rotation up to $30^\circ$, crops of $55$--$100\%$, brightness/contrast/saturation factors in $[0.6,1.4]$, and up to two of grayscale, blur, sharpen, equalise, autocontrast, posterise, edge detection and inversion), and embed the originals and copies. With $E_0$ the original embeddings, $\Sigma_b$ is the covariance of $E_0$ and $\Sigma_w$ the covariance of the differences between copies of one image (both copies against the original and against each other, halved). We solve $\Sigma_b v=\lambda(\Sigma_w+\epsilon I)v$ with $\epsilon=10^{-3}\,\mathrm{tr}(\Sigma_w)/d$ and keep the top $128$ directions. The edit family is fixed before seeing any release and is the same for all datasets.

\section{Additional results}
\label{app:orig}
\label{app:backbones}

\input{tables/real.tex}
\input{tables/exposure.tex}
\input{tables/orig.tex}
\input{tables/datasize.tex}
\input{tables/backbones.tex}
\input{tables/resid.tex}
\input{tables/variants.tex}

\input{tables/audit.tex}

\FloatBarrier
\section{Fine-tuning}
\label{app:finetune}
\IfFileExists{tables/finetune.tex}{\input{tables/finetune.tex}}{Fine-tuning results are pending.}
